% !TEX program = pdflatex
\documentclass[11pt]{article}
\usepackage[T1]{fontenc}
\usepackage{lmodern,amsmath,amssymb,amsthm,microtype,booktabs}
\usepackage[margin=1in]{geometry}
\usepackage[hidelinks]{hyperref}
\hypersetup{pdftitle={Operator pairs and all bounded profiles of a six-dimensional dg algebra}}
\newtheorem{theorem}{Theorem}
\newtheorem{lemma}[theorem]{Lemma}
\newtheorem{corollary}[theorem]{Corollary}
\newcommand{\per}{\operatorname{per}}
\newcommand{\fd}{\mathrm{fd}}
\newcommand{\HHom}{\operatorname{RHom}}
\newcommand{\Hom}{\operatorname{Hom}}
\newcommand{\Ext}{\operatorname{Ext}}
\newcommand{\cone}{\operatorname{Cone}}
\newcommand{\thick}{\operatorname{thick}}
\newcommand{\supp}{\operatorname{supp}_{\sigma}}
\newcommand{\Dom}{\operatorname{Dom}_{\sigma}}
\newcommand{\prof}{\mathcal U}
\newcommand{\fib}[1]{\overline{#1}}
\newcommand{\bb}{\mathcal B}
\newcommand{\ttt}{\mathcal T}
\setlength{\emergencystretch}{1em}
\title{Operator pairs and all bounded profiles\\ of a six-dimensional dg algebra}
\author{Marouane Lamharzi Alaoui\\
\small Research notes on two-term dg kernels and bounded profiles of universal localisations}
\date{18 November 2026}
\begin{document}
\maketitle
\begin{abstract}
Two finite-dimensional operator pairs give the same bounded profile precisely
when their sets of simple composition factors agree after identifying a
factor with its exchange. We describe the whole poset over any field.
One part consists of finite sets of these simple types in reverse inclusion.
The other consists of invertibility domains of square matrices over a
two-generator free algebra, tested in both orders of the generators.
A special-fibre calculation separates the two parts. We give a finite
comparison method for operator pairs and for polynomial-matrix domains.
Over $\mathbb C$ we identify the
nonempty domains by finite sets of stable atom types.
\end{abstract}

\section{The classification}
Let $D$ be the six-dimensional dg algebra over $k$ with degree-zero
basis $1,a,b,h$, degree-minus-one basis $r,w$, products $a^2=h$,
$br=rb=w$, and $dr=h$. All unlisted radical products and differentials
vanish. Its finite-cohomology category is equivalent to the perfect
\emph{right} dg modules over the zero-differential graded algebra
\[
 A=k\langle x,z\rangle/(x^2z-zx^2),\qquad |x|=|z|=1 .
\]
The preceding note \emph{A cubic higher multiplication creates infinitely
many bounded profiles} proves this using the reduced bar dual. We work
on the $A$ side. For perfect $P,E$, membership $P\in\prof_E$ means that
$H^*\HHom_A(P,E)$ has finite total dimension. Such Hom complexes have
locally finite, bounded-below cohomology, by building both objects from
finite cells and retracts. The inclusion poset is the same after dividing
by the original $\per D$, common to all profiles.

Put $s=x^2$, $R=k\langle U,V\rangle$, and let $\sigma$ exchange $U,V$.
Let $\Omega_r$ and $\Omega_l$ be the isomorphism types of nonzero
finite-dimensional simple right and left $R$-modules, respectively,
modulo $\sigma$. For a finite module $W$ write $\supp W$ for its
set of such orbits among its composition factors, with repetitions
discarded. A superscript $\sigma$ denotes twisting the action. Ordinary
matrices on a vector space specify either convention using rows or
columns consistently. Denote by $M(W)$ the perfect module of the preceding
note, with $W$ in every nonnegative degree, $x$ the identity between
consecutive copies and $z$ the two operators alternately. Its left-handed
version is perfect too, by word reversal.

For a square polynomial matrix $F$ over $R$, set
\[
 \Dom F=\{[L]\in\Omega_l:\ F(L)\text{ and }F(L^\sigma)
                      \text{ are invertible}\}.
\]
Here $F(L)$ is the block matrix acting on copies of $L$. Include the
empty domain (zero matrix) and all of $\Omega_l$ (identity matrix).
The use of orbits puts the parity condition into the label.

\begin{theorem}\label{poset}
The list below has no repetitions except for equal labels within a row.
Every label in the list occurs.
\begin{center}
\begin{tabular}{lll}
\toprule
Profiles & Labels & Inclusion rule\\
\midrule
$\ttt_F$ & finite $F\subset\Omega_r$ & $\ttt_F\subseteq\ttt_{F'}\iff F\supseteq F'$\\
$\bb_O$ & domains $O=\Dom H\subseteq\Omega_l$ &
$\bb_O\subseteq\bb_{O'}\iff O\subseteq O'$\\
\bottomrule
\end{tabular}
\end{center}
Here $\prof_{M(W)}=\ttt_{\supp W}$. The full profile is $\ttt_\varnothing$,
and the smallest is $\bb_\varnothing$. A first-row profile never lies
in a second-row profile as a subcategory. A second-row profile lies
in a first-row profile exactly when the former is the smallest or the
latter is the full one. Thus every nonempty $F$ and nonempty $O$ give
incomparable profiles with intersection $\bb_\varnothing$.
\end{theorem}

In particular a Jordan block can have the same profile as a single scalar,
and the exchange can be made independently on different simple factors.
This classification by actual subsets of simple types makes sense over
any field. Sections~\ref{atoms} and~\ref{effective} give finite-factor labels over
$\mathbb C$ and a finite comparison of two supplied domains over any field.

\section{The localized presentation}
Write $B=A[s^{-1}]$. The normal words for $A$ avoid $xxz$.
Right multiplication by $x^2$ appends two letters to these words and is
injective. Since $s$ is central, its localisation is flat and
$G=A/(s)=k\langle x,z\rangle/(x^2)$. There is a graded isomorphism
\[
 B=R[x,x^{-1};\sigma],\qquad U=zx^{-1},\quad V=xzx^{-2},
 \qquad x f=\sigma(f)x .
\]
Degree zero and induction from $R=B^0$ are inverse equivalences on
graded modules; multiplying by the invertible $x$ identifies the
components, with $\sigma$ twisting the action after one step.

\begin{lemma}\label{presentation}
Any perfect right dg $B$-module is isomorphic in the derived category to
the cone of a map $B^m\to B^n$ with entries in $R$ and generators in
degree zero. If its degree-zero cohomology is a finite-dimensional module
$W$, its cohomology as a graded module is the two-sided infinite grading
of $M(W)$, and it is quasi-isomorphic to that cohomology with zero
differential. The corresponding statements hold for left modules.
\end{lemma}
\begin{proof}
In the category of graded $B$-modules, cohomology of a perfect dg module
is finitely presented. Indeed $R$ is coherent and hereditary with free
submodules of a free module. For the freeness assertion take leading
terms in a degree-lexicographic order on module words, compatible with
right multiplication. Pick a monic element for each minimal prefix among
the leading words in each position. Their right multiples span by
leading-term reduction and are independent by prefix incomparability.
A finitely generated free submodule has finite rank. The image of a
map of finite free modules is thus finite free, and its kernel is a
finitely generated projective summand, also free. This gives coherence.
The graded equivalence transfers these facts to $B$.
The long cohomology sequence now shows that finite presentation is closed
under triangles (kernels, cokernels and extensions), shifts and summands.

Take a finite graded free presentation $0\to B^m\xrightarrow{F} B^n
\to H^*P\to0$, using the equivalence with $R$ to put the generators
in degree zero. Lift the generators of $B^n$ to cycles in $P$.
The relations are boundaries, so lift each to a primitive. The resulting
map $\cone(F)\to P$ induces the presentation isomorphism on cohomology.
Also $\cone(F)\to H^*P$ is a quasi-isomorphism by injectivity of $F$.
Degree zero determines the graded module, giving the final claim.
\end{proof}

For $E$ perfect call its cohomology \emph{eventually $s$-invertible}
if $s:H^nE\to H^{n+2}E$ is an isomorphism for all sufficiently large
$n$. Its localisation map to $E_s$ is then an isomorphism on all high
cohomology groups. The same terminology applies to $P^\vee=\HHom_A(P,A)$,
a perfect left module. If $W=H^0(E_s)$ in this case, $W$ is finite
dimensional. For the alternating detector $M(W)$ it is exactly the
original ordered pair, since $zx^{-1}$ and $xzx^{-2}$ act on degree zero
as the even and odd operators. The conjugate occurs in odd degrees.

\section{Equality and inclusion for operator pairs}
For perfect $P$, the localization of $\HHom_A(P,E)$ is
$\HHom_B(P_s,E_s)$. This follows first for $P=A$, then for its thick
closure. If $E$ has eventually invertible $s$, Hom into the cone of
$E\to E_s$ vanishes in high degrees for every perfect $P$, again by
finite cells and summands. The localized Hom has two-periodic cohomology.
Consequently $P\in\prof_E$ exactly when $\HHom_B(P_s,E_s)$ is acyclic.
Use Lemma~\ref{presentation} to present $P_s$ by a matrix $H$.
Applying Hom into the zero-differential model of $E_s$ tests bijectivity
of its precomposition map on $W$ and $W^\sigma$. For $W\ne0$ unequal
matrix dimensions preclude bijectivity. For a square matrix, a composition
series puts each evaluation into block triangular form with square
diagonal blocks. Bijectivity is therefore equivalent to bijectivity on
each simple factor in both parities. This proves sufficiency of the
support and reverse inclusion, also for sources which are retracts.

For necessity take a simple type present in $W'$ and absent up to exchange
from $W$. The annihilator of a finite-dimensional simple module is a
maximal two-sided ideal: the image on the simple is a finite-dimensional
primitive algebra, hence simple artinian. Equal annihilators give
isomorphic simples on the same side. In the semisimple quotient of the
joint finite image on all the relevant simple representatives and
exchanges, choose the central idempotent equal to one on the support
of $W$ and zero on the distinguished orbit. Equivalently lift the
prescribed scalars by the Chinese remainder theorem for distinct
maximal two-sided ideals. This yields $f\in R$ with both properties.
As $f\in B^0$, some $s^d f\in A^{2d}$. The cone of left multiplication
by $s^d f:A[-2d]\to A$ gives a perfect source in $\prof_{M(W)}$
outside $\prof_{M(W')}$. This works if $W=0$ as well. Taking direct sums
of simple representatives realises every finite support.

\paragraph{A finite comparison on the input matrices.}
Form the unital joint image algebra $H$ on $W\oplus W^\sigma\oplus
W'\oplus(W')^\sigma$ by closing the span of the identity under the two
generators. Let $I$ be the kernel of restriction to $W\oplus W^\sigma$.
Then $\prof_{M(W)}\subseteq\prof_{M(W')}$ exactly when repeated action
of $I$ on $W'\oplus(W')^\sigma$ reaches zero. Indeed a faithful finite
module contains every simple type of its algebra among its composition
factors: the regular module embeds in finitely many copies by its action
on a vector-space basis. Apply this to $H/I$. On a simple module the
ideal $I$ acts either by zero or surjectively. A finite filtration now
proves the criterion in both directions. Each subspace iteration
decreases dimension or stabilises, so this test terminates over the
given field using only linear algebra. Interchanging the inputs tests
the reverse inclusion.

\section{Why the special fibre forces the other cases}
For a perfect complex use a subscript $s$ for central localization and
a bar for derived reduction to $G$. Let $K=\HHom_A(P,E)$. Flat localization
and the two-term free resolution of $G$ give
\[
 K_s\simeq E_s\otimes^L_B (P^\vee)_s,\qquad
 \cone(s:K[-2]\to K)\simeq\HHom_G(\fib P,\fib E).
\]
The second identity also follows by applying $\HHom_A(P,-)$ to the
multiplication triangle for $E$ and using extension-restriction
adjunction. Reduction commutes with the perfect dual.
For any locally finite bounded-below cohomology as here, $H^*K$ is
finite exactly when $K_s$ is acyclic and the displayed cone has finite
cohomology. The cone condition makes $s$ eventually invertible, and
localization then recovers each high-degree group.

\begin{lemma}\label{fibre}
For perfect right dg $G$-modules $P_0,E_0$, the total cohomology of
$\HHom_G(P_0,E_0)$ is finite exactly when
$H^*(E_0)$ is finite or $H^*\HHom_G(P_0,G)$ is finite.
\end{lemma}
\begin{proof}
Use the ordinary commutative local ring
$T=k[a,b]_{(a,b)}/(ab,b^2)$ with residue field $k$. We verify the
Koszul duality needed here directly. A minimal free $T$-resolution
$J$ of $k$ has basis $e_w$ in degree $-|w|$ indexed by words on $a,b$
avoiding $aa$ and differential $d(e_{iw})=i e_w$.
The initial augmentation is the residue map. For each suffix the images
$aT,bT$ intersect trivially, $\operatorname{ann}(a)=kb$ and
$\operatorname{ann}(b)=(a,b)$. These facts prove exactness recursively,
including at suffixes beginning in $a$ where only the prefix $b$ is allowed.
There are two closed degree-one endomorphisms stripping a terminal $a$
or $b$ with sign $(-1)^{|w|-1}$ on the entire input word $w$ (zero on
the empty word or a mismatched ending). Stripping $a$ twice is zero.
Products for admissible words evaluate, up to signs, as the dual basis
of $\Hom_T(J,k)$. Since $J$ is bounded-above projective,
$\operatorname{End}_T(J)\to\Hom_T(J,k)$ induced by augmentation is a
quasi-isomorphism of complexes. Thus $G\to\operatorname{End}_T(J)$ is
a dg algebra quasi-isomorphism. The resulting equivalence
$\thick_{D(T)}(k)\simeq\per G$ takes $k$ to $G$. Its source is precisely
the bounded complexes with finite-length cohomology, by truncation.

Write $Y,N$ for the objects corresponding to $P_0,E_0$. If
$\Ext_T^*(k,N)$ is bounded, construction from $k$ makes
$\Ext_T^*(Y,N)$ bounded for all such $Y$. Otherwise take a bounded-above
minimal free resolution of $Y$ with finite ranks. Below its cohomology,
dimension shifting replaces $Y$, up to a perfect complex and a shift,
by a syzygy $L\subseteq(a,b)T^n$ (truncate the resolution in an exact
degree). If $Y$ is nonperfect this syzygy is nonzero. It is killed by
$b$ and is a submodule of $a k[a]_{(a)}^n\oplus (kb)^n$. Over the
discrete valuation ring $T/(b)$, its torsion is killed by $a$, so
$L\cong (T/(b))^u\oplus k^v$ with $u+v>0$. The first syzygy of $T/(b)$
is $kb$ and that of $k$ is $aT\oplus kb$. Thus the next syzygy of $L$
has a $k$ summand. Unbounded $\Ext_T^*(k,N)$ now forces unbounded
$\Ext_T^*(Y,N)$. Bounded total Ext from $Y$ to $k$ is equivalent to
perfection of $Y$ by minimality. Conversely a perfect $Y$ has bounded
Ext to every $N$ here. Translating the two tests back gives the assertion.
\end{proof}

Combining the lemma and the multiplication triangles yields an individual
membership test, useful even without comparing poset labels:
\[
 \boxed{\quad P\in\prof_E\ \Longleftrightarrow E_s\otimes_B^L(P^\vee)_s\simeq0\quad\text{and}\quad
 \big(E\text{ or }P^\vee\text{ is eventually }s\text{-invertible}\big).\quad}
\]
Here and below eventual invertibility refers to cohomology. In particular
the perfect source $G=\cone(s:A[-2]\to A)$ belongs to $\prof_E$
exactly when $E$ is eventually $s$-invertible. Such profiles are exactly
the first row of Theorem~\ref{poset}.

\section{Matrix domains and the order between the two parts}
Suppose $E$ fails eventual invertibility. By the boxed criterion, the
eligible sources have $P^\vee$ eventually $s$-invertible. Put
$W_P=H^0((P^\vee)_s)$, a finite-dimensional left $R$-module.
Lemma~\ref{presentation} gives a presentation $B^m\xrightarrow{H}B^n$
for $E_s$. Tensor with the zero-differential localized model of $P^\vee$.
In each parity this evaluates $H$ on copies of $W_P$ or $W_P^\sigma$.
If $m\ne n$, acyclicity requires $W_P=0$. Otherwise the composition-series
test gives the explicit rule
\[
  P\in\bb_O\quad\Longleftrightarrow\quad P^\vee\text{ is eventually }
  s\text{-invertible and }\supp W_P\subseteq O,
  \qquad O=\Dom H .
\]
The unequal-rank case is $O=\varnothing$. Any finite-dimensional left
module can occur as $W_P$, by taking $P$ to be the perfect dual of its
left alternating detector. This both witnesses differences of domains
and proves their inclusion rule.

Conversely given any square matrix $H$ over $R$, clear all central
denominators by a common $s^d$. Let $E$ be the direct sum of $G$ and
the cone of the resulting graded map $A[-2d]^m\to A^m$. Its $G$
summand has infinite cohomology killed by $s$, and its localization is
zero. The other summand realizes exactly the desired evaluation test
after localization. This argument applies even when the matrix over $R$
has a kernel. In particular $\bb_{\Omega_l}=\prof_G$ and
$\bb_\varnothing=\prof_A$. The latter consists of the sources whose
dual has finite total cohomology. Under the original equivalence
$A$ is the residue-field detector, so this is the image of $\per D$.

To check the mixed comparisons, the source $G$ lies in every first-row
profile and no second-row profile. Take $P$ in a second-row profile and
a nonzero $W$ giving a first-row detector. If $W_P\ne0$, the localized
tensor cannot be acyclic. Indeed, for a finite-dimensional right
$R$-module $W$ of dimension $d$, there is a free presentation with
ranks $2d,d$. Start with $W\otimes_k R\to W$ and impose the relations
$v\otimes i-(vi)\otimes1$ for $i=U,V$. Their right multiples are
independent by the leading first letters and they span the kernel by
word reduction. Induce this presentation to $B$ and tensor with
$(P^\vee)_s$. The dimensions in each parity differ by $d\dim W_P>0$.
Hence intersection with the first-row profile forces $W_P=0$, which
here is equivalent to $P\in\bb_\varnothing$. Conversely all such
sources pass every test. This proves Theorem~\ref{poset}.

\section{Counterexamples and the complex case}\label{atoms}
\paragraph{Separate similarity and extensions.}
Let $J=e_{12}$ and $K=e_{21}$ on $k^2$. The pairs $(J,K)$ and $(J,J)$
have individually similar first operators and individually similar
second operators. The polynomial $U-V$ is invertible on the former
and zero on the latter, also after exchange, over every field.
Clearing its denominator gives the explicit cone of $zx-xz$.
The pair $(J,J)$, although nonsplit, has exactly the same profile as
the one-dimensional zero pair. Both have only the trivial simple type.
Over a field containing $0,1,2$ as distinct elements, the sums
$(0,1)\oplus(0,2)$ and $(1,0)\oplus(0,2)$ have equal profiles, although
neither simultaneous similarity nor one global exchange identifies them.

\paragraph{Even the entire scalar pencil can agree.}
Over $\mathbb Q$ take $L=\operatorname{diag}(2,3,4)$ and let $S$ cyclically
send $e_1$ to $e_2$, $e_2$ to $e_3$, and $e_3$ to $e_1$. Compare
$(L,S)$ with $(L,S^{\mathsf t})$. Their scalar pencils
$\alpha I+\beta L+\gamma S$ and $\alpha I+\beta L+\gamma S^{\mathsf t}$
have identical determinant polynomials by transposition. Set
\[
 f(U,V)=[U,V]^3-2I+(U-2I)(U-3I)(U-4I).
\]
Since $[L,S]^3=-2I$, the two evaluations on the first pair and its
exchange are $-4I$ and $(S-2I)(S-3I)(S-4I)$, both invertible.
On $(L,S^{\mathsf t})$ the commutator cube is $2I$ and $f$ is zero.
Thus clearing denominators of this single polynomial separates the
profiles. The pairs are both simple: the spectral projections of $L$
and the directed cycle generate the matrix units. Their distinct swap
orbits give incomparability in both directions.

\paragraph{The localization can forget the smallest profile.}
For $W\ne0$, $M(W)$ and $M(W)\oplus G$ have identical localizations.
The first yields $\ttt_{\supp W}$, whereas the second yields
$\ttt_{\supp W}\cap\bb_{\Omega_l}=\bb_\varnothing$.
The special-fibre condition in the membership test is essential.

\paragraph{Scalar tests miss second-row profiles.}
The commutator $H=UV-VU$ vanishes on all one-dimensional modules. It
evaluates on $(J,K)$ as $\operatorname{diag}(1,-1)$, invertible also in
characteristic two and after exchange. The pair $(J,K)$ is simple
because $J,K$ generate all matrix units. Thus $\bb_{\Dom H}$ is
strictly above $\bb_\varnothing$ and strictly below $\bb_{\Omega_l}$.
A source witnessing the first strict inclusion is the dual alternating
detector for $(J,K)$; the one-dimensional zero pair witnesses the second.
This example also gives a nontrivial second-row profile incomparable
with every proper first-row profile.

\paragraph{Finite factor labels over $\mathbb C$.}
A square matrix over $\mathbb C\langle U,V\rangle$ is full if it has
full inner rank. Fullness is equivalent to admitting an invertible
finite matrix evaluation (the free skew field criterion). A full
nonunit has a factorisation into atoms. Stable association allows
adjoining identity blocks and multiplying on both sides by invertible
polynomial matrices. For two full matrices $H,H'$, use the
Singul\"arstellensatz, Theorem 2.12 of
\href{https://arxiv.org/abs/1907.04328}{Helton, Klep and Vol\v ci\v c}.
The free singularity locus of $H$ is contained in that of $H'$ precisely
when every atomic factor of $H$ occurs, up to stable association, in $H'$.
Passing from simple tuples to all finite tuples preserves comparisons of
invertibility domains by block triangularity. Testing both parities is
equivalent to the single block $H\oplus\sigma H$. Consequently the
nonempty second-row domains are in bijection with finite sets of
stable atom types modulo $\sigma$, ordered by reverse inclusion.
Every finite set occurs by taking block sums. A block sum of full matrices
is full over the free skew field, so the symmetrization still has a
nonempty domain. Nonfull matrices give the global bottom, and the
identity has empty atom set and gives $\prof_G$.
The arbitrary-field description in Theorem~\ref{poset} uses no
characteristic-zero factorisation theorem.

\section{Comparing the second-row labels over any field}\label{effective}
The description by domains has a finite pairwise comparison test as well.
This section gives it for two given polynomial-matrix labels; it is
separate from the much smaller calculation for input operator pairs.

\begin{lemma}\label{monic}
For monic pencils $I-\sum_i C_iY_i$ and $I-\sum_i C'_iY_i$, the first
invertibility domain on finite matrix tuples is contained in the second
precisely when the nontrivial simple factors of the coefficient tuple $C'$
are contained among those of $C$, ignoring multiplicities. Here the
trivial simple has every generator zero. This assertion holds over any field.
\end{lemma}
\begin{proof}
Triangularising along a composition series gives sufficiency; the trivial
factor contributes an identity. If a nontrivial factor of $C'$ is missing,
use the finite-image argument of Section 3 to find a polynomial $f$ in
the coefficient variables which is one on that factor and zero on all
factors of $C$ and on the trivial module. In particular $f$ has zero
constant term. There exist finite scalar matrices $Y_i$ that make the
pencil invertible precisely when $I-f(C)$ is invertible, for every $C$.
To construct them write $f$ as a sum of words. For each word make a
path from one common base vertex back to itself, with distinct internal
vertices, edge labels the letters of the word and its coefficient on one
edge. Let $Y_i$ be the weighted adjacency matrices for label $i$, adding
parallel edge weights. Deleting the base vertex leaves a directed acyclic
graph. Eliminating its vertices in the block matrix $I-\sum_iY_i\otimes C_i$
has identity diagonal pivots and sole remaining Schur complement $I-f(C)$.
Interchanging tensor factors is a change of basis. This gives a tuple
where the first pencil is invertible and the second singular.
\end{proof}

\begin{corollary}\label{finite}
Inclusion and equality between two second-row labels can be checked by
finitely many field operations and polynomial identity tests, in any
characteristic, given exact arithmetic in $k$.
\end{corollary}
\begin{proof}
First replace each label $H$ by $H\oplus\sigma H$. Its ordinary domain
on all finite tuples determines the original domain on simple orbits
and conversely, by composition series. Linearise each matrix to an affine
pencil $L_i=A_{i0}+A_{i1}U+A_{i2}V$ preserving invertibility at every tuple.
For example replace a product term $BC$ in a matrix $D+BC$ by the block
matrix with rows $(D,B),(-C,I)$ and use its Schur complement. Splitting
individual word products, with the appropriate entry row and column,
terminates in affine entries.

Let $n_i$ be the pencil sizes and $d=\max(1,n_1-1,n_2-1)$. Test the
determinants $\det L_i(Z_1,Z_2)$ at two $d$-by-$d$ independent generic
matrices. An identically zero determinant means an empty domain.
Indeed apply Theorem 1.8 of
\href{https://arxiv.org/abs/1512.03393}{Derksen and Makam} over the
algebraic closure to the homogeneous pencil
$\sum_{j=0}^2 A_{ij}\otimes Z_j$. Existence of any invertible evaluation
implies one in every size $d\ge n_i-1$ for $n_i\ge2$; the scalar
case is immediate. In fixed size the open invertibility condition can
be combined with $Z_0$ invertible. Multiplying by $I\otimes Z_0^{-1}$
then normalizes $Z_0$ to identity. Conversely a nonzero generic determinant
can be specialized in the algebraic closure, and restricting scalars from
the finite extension generated by the entries produces a tuple over $k$.
This also justifies the emptiness test over finite fields. Empty-domain
comparisons follow directly.

For two nonempty domains take the generic pair $Z$ as a common centre
over $K=k(\text{entries of }Z)$. Inclusion of domains on all finite tuples
over $k$ is equivalent to inclusion over $K$. A witness over $k$ persists.
A witness over $K$ can be specialized to a finite algebraic extension of
$k$, avoiding denominators and the nonzero determinant, while its identically
zero determinant remains zero. Restriction of scalars preserves the two
outcomes. At sizes $dm$ now substitute
\[
 Z_i\otimes I_m+\sum_{a,b=1}^d e_{ab}\otimes Y_{iab}
 \quad(i=1,2)
\]
with the $2d^2$ blocks free to vary. This covers every tuple of size $dm$
over $K$. Any witness in another size can be amplified by a direct sum of
$d$ copies. The constant part of the amplified affine pencil is invertible
over $K$; left normalization gives a monic pencil in the block variables.
Lemma~\ref{monic} compares its ordinary domains. To implement that finite
simple-factor test, adjoin a trivial summand to both coefficient tuples
and use the joint-image ideal iteration of Section 3, with all $2d^2$
generators and without any further exchange. All steps use finite matrices
and rational functions. This is a termination statement, without a
practical complexity estimate.
\end{proof}

The preceding factor theorem and the monic-pencil comparison of
\href{https://arxiv.org/abs/1512.02648}{Klep and Vol\v ci\v c}
supply prior free-locus results. The special-fibre argument is in the
setting of the vanishing theorem for nontrivial fibre products of
\href{https://doi.org/10.1090/proc/13633}{Nasseh and Sather-Wagstaff}:
$T$ is the fibre product of a discrete valuation ring and the dual
numbers over $k$. The proof above establishes the particular complex
version directly. The classification here combines this fibre, the
graded localization and the realisation of all test matrices for the
six-dimensional example. Its order description does not classify
simultaneous similarity classes of all operator pairs.
\end{document}
